\documentclass[11pt]{article}
\usepackage[a4paper,margin=1in]{geometry}
\usepackage{amsmath,amssymb}
\usepackage{enumitem}

\begin{document}

\begin{center}
\textbf{CSE 400: Fundamentals of Probability in Computing} \\
\textbf{Lecture 13 (Tutorial) — Exam-Oriented Scribe}
\end{center}

Based strictly on Tutorial-1 (February 13, 2026) 

Tutorial\_1

\section*{\textbf{Q1. Division of 20 Distinct Dishes into Four Platters}}

\textbf{Problem Data}

Total dishes: 20 (all distinct).

To be divided into 4 platters.

Each platter must contain exactly 5 dishes.

The order of platters does not matter.

\subsection*{\textbf{(a) Total Number of Ways}}

\textbf{Definitions and Assumptions}

We partition 20 distinct objects into 4 subsets of equal size (5 each).

Subsets are unordered.

\textbf{Step 1: Count Ordered Selections}

Choose 5 dishes for the first platter:

\[
\binom{20}{5}
\]

Choose 5 of the remaining 15:

\[
\binom{15}{5}
\]

Choose 5 of the remaining 10:

\[
\binom{10}{5}
\]

Choose the remaining 5:

\[
\binom{5}{5}
\]

Thus, total ordered groupings:

\[
\binom{20}{5}
\binom{15}{5}
\binom{10}{5}
\binom{5}{5}
\]

\textbf{Step 2: Remove Ordering of Platters}

Since platters are indistinguishable, divide by $4!$:

\[
\frac{1}{4!}
\binom{20}{5}
\binom{15}{5}
\binom{10}{5}
\binom{5}{5}
\]

\subsection*{\textbf{(b) Probability All Platters Contain One Cuisine Each}}

\textbf{Additional Information}

5 dishes from each of four cuisines.

Total dishes = 20.

\textbf{Total Possible Divisions}

Same as in part (a):

\[
\frac{1}{4!}
\binom{20}{5}
\binom{15}{5}
\binom{10}{5}
\binom{5}{5}
\]

\textbf{Favorable Divisions}

Each cuisine has exactly 5 dishes.

Hence each cuisine must form one platter.

Only the assignment of cuisines to platters varies.

Number of assignments:

\[
4!
\]

\textbf{Probability}

\[
\frac{4!}{
\frac{1}{4!}
\binom{20}{5}
\binom{15}{5}
\binom{10}{5}
\binom{5}{5}}
\]

\section*{\textbf{Q2. Waiting Time with Uniform Arrival}}

\textbf{Problem Data}

Buses arrive every 15 minutes starting at 7:00 A.M.

Passenger arrival time uniformly distributed between 7:00 and 7:30.

Let arrival time $T \sim U(0,30)$.

\subsection*{\textbf{(a) Waiting Less Than 5 Minutes}}

Waiting time is the time until the next bus.

Bus arrival times:

0 minutes

15 minutes

30 minutes

Arrival results in waiting < 5 minutes if:

Arrival between 10 and 15,

Arrival between 25 and 30.

Total favorable length:

\[
5+5=10
\]

Total interval length:

30

\[
P(wait<5)=\frac{10}{30}=\frac{1}{3}
\]

\subsection*{\textbf{(b) Waiting More Than 10 Minutes}}

Waiting > 10 minutes if:

Arrival between 0 and 5,

Arrival between 15 and 20.

Total favorable length:

\[
5+5=10
\]

\[
P(wait>10)=\frac{10}{30}=\frac{1}{3}
\]

\section*{\textbf{Q3. Conditional Probability}}

Given

\[
P(L)=0.15
\]

\[
P(E)=0.25
\]

\[
P(L \cap E)=0.08
\]

Where:

$L$: arrives late,

$E$: leaves early.

We seek:

\[
P(L^c \mid E^c)
\]

\textbf{Step 1: Complements}

\[
P(L^c)=1-0.15=0.85
\]

\[
P(E^c)=1-0.25=0.75
\]

\textbf{Step 2: Inclusion–Exclusion}

\[
P(L \cup E)=P(L)+P(E)-P(L \cap E)=0.15+0.25-0.08=0.32
\]

\textbf{Step 3: Complement Event}

\[
P(L^c \cap E^c)=1-P(L \cup E)=1-0.32=0.68
\]

\textbf{Step 4: Conditional Probability}

\[
P(L^c \mid E^c)=\frac{P(L^c \cap E^c)}{P(E^c)}=\frac{0.68}{0.75}
\]

\section*{\textbf{Q4. Poisson Distribution (Mean 0.2)}}

\textbf{Assumption}

Expected number of typographical errors per page = 0.2.

$X \sim Poisson(0.2)$

PMF:

\[
P(X=k)=e^{-0.2}\frac{0.2^k}{k!}
\]

\subsection*{\textbf{(a) Zero Errors}}

\[
P(X=0)=e^{-0.2}
\]

\subsection*{\textbf{(b) Two or More Errors}}

\[
P(X \ge 2)=1-P(0)-P(1)
\]

\[
P(1)=e^{-0.2}(0.2)
\]

\[
P(X \ge 2)=1-e^{-0.2}(1+0.2)
\]

\section*{\textbf{Q5. Poisson Distribution (Mean 3.5)}}

$X \sim Poisson(3.5)$

\subsection*{\textbf{(a) At Least 2 Accidents}}

\[
P(X \ge 2)=1-P(0)-P(1)=1-e^{-3.5}(1+3.5)
\]

\subsection*{\textbf{(b) At Most 1 Accident}}

\[
P(X \le 1)=e^{-3.5}(1+3.5)
\]

\section*{\textbf{Q6. Poisson ($\lambda = 3$)}}

$X \sim Poisson(3)$

\subsection*{\textbf{(a) $P(X \ge 3)$}}

\[
P(X \ge 3)=1-P(0)-P(1)-P(2)
\]

\[
=1-e^{-3}\left(1+3+\frac{9}{2}\right)
\]

\subsection*{\textbf{(b) Conditional Probability}}

\[
P(X \ge 3 \mid X \ge 1)=\frac{P(X \ge 3)}{P(X \ge 1)}
\]

\[
P(X \ge 1)=1-e^{-3}
\]

\section*{\textbf{Q7. Continuous Random Variable}}

Given PDF

\[
f_X(x)=\frac{1}{8\pi}\exp\left(-\frac{(x+3)^2}{8}\right)
\]

Mean = $-3$

Variance = $4$

Standard deviation = $2$

Standardization:

\[
Z=\frac{X+3}{2}
\]

1. 
\[
P(X \le 0)
\]

\[
Z=\frac{0+3}{2}=1.5
\]

\[
P(X \le 0)=1-Q(1.5)
\]

2. 
\[
P(X>4)
\]

\[
Z=\frac{4+3}{2}=3.5
\]

\[
P(X>4)=Q(3.5)
\]

3. 
\[
P(|X+3|<2)
\]

\[
|Z|<1
\]

\[
=1-2Q(1)
\]

4. 
\[
P(|X-2|>1)
\]

\[
P(X>3)+P(X<1)
\]

Express in Q-function form via standardization.

\section*{\textbf{Q8. Jury Decision}}

Given

12 jurors.

At least 8 must vote guilty.

Each juror independently makes the correct decision with probability $\theta$.

Let:

$X \sim Binomial(12,\theta)$

Correct decision if:

\[
X \ge 8
\]

\[
P(Correct)=\sum_{k=8}^{12} \binom{12}{k}\theta^k(1-\theta)^{12-k}
\]

\section*{\textbf{Q9. PMF}}

\[
p(i)=c\frac{\lambda^i}{i!}
\]

Determination of $c$

\[
\sum_{i=0}^{\infty} c\frac{\lambda^i}{i!}=1
\]

\[
ce^{\lambda}=1
\]

\[
c=e^{-\lambda}
\]

\subsection*{\textbf{(a)}}

\[
P(X=0)=e^{-\lambda}
\]

\subsection*{\textbf{(b)}}

\[
P(X>2)=1-e^{-\lambda}\left(1+\lambda+\frac{\lambda^2}{2}\right)
\]

\section*{\textbf{Q10. System Reliability}}

Given

$n$ independent components.

Each functions with probability $p$.

System operates if at least half function.

Let:

$X \sim Binomial(n,p)$

\subsection*{\textbf{(a) 5 vs 3 Components}}

5-component works if $X \ge 3$:

\[
P_5=\sum_{k=3}^{5}\binom{5}{k}p^k(1-p)^{5-k}
\]

3-component works if $X \ge 2$:

\[
P_3=\sum_{k=2}^{3}\binom{3}{k}p^k(1-p)^{3-k}
\]

Compare:

\[
P_5>P_3
\]

\subsection*{\textbf{(b) General Case}}

Compare:

\[
\sum_{k=k+1}^{2k+1}\binom{2k+1}{k}p^k(1-p)^k
\]

with

\[
\sum_{k=k}^{2k-1}\binom{2k-1}{k}p^k(1-p)^{k-1}
\]

This determines when the $(2k+1)$-component system is better than the $(2k-1)$-component system.

End of Lecture 13 Tutorial Scribe

\end{document}